\section{Level-matched short excursions}
\label{sec:v74-level-matching}
The two Jordan masses near a roof need not describe the same
excursion. We pair entrances and exits on the same component of a
superlevel set. This construction is made after summing the entire
physical density, rather than separately on words or threshold sources.

Let $f\ge0$ be compactly supported and in $BV(\R)$. For almost every
$u>0$, the set $E_u=\{f>u\}$ agrees modulo null sets with a finite
disjoint union of bounded open intervals. Use that essential interval
representative. If $t$ lies in its component $I_u(t)$, write
$\ell_u(t)=|I_u(t)|$. For almost every $(u,t)$ with $0<u<f(t)$,
this length is finite and positive. Define
\begin{equation}\label{eq:v74-short-excursion}
 S_\delta f(t)=\int_0^{f(t)}
              \left(1-\frac{\ell_u(t)}{\delta}\right)_+du,
 \qquad \delta>0.
\end{equation}
Exceptional levels and pairs are assigned zero. All identities and
norms below are almost-everywhere identities and essential norms.

\begin{lemma}[Measurability and oriented level decomposition]
\label{lem:v74-level-coarea}
The function in \eqref{eq:v74-short-excursion} is measurable. As
finite measures,
\begin{equation}\label{eq:v74-oriented-levels}
 (Df)^+=\int_0^\infty\sum_{(a,b)\subset E_u}\delta_a\,du,
 \qquad
 (Df)^-=\int_0^\infty\sum_{(a,b)\subset E_u}\delta_b\,du.
\end{equation}
In each sum, $(a,b)$ runs over the essential interval components.
\end{lemma}
\begin{proof}
The one-dimensional $BV$ coarea theorem gives finite perimeter of
$E_u$ for almost every $u$, and the finite-interval description follows
since these sets are bounded. Endpoints can be chosen measurably:
for an interior pair $(u,t)$ take the infimum of rational $a<t$
for which $\int_a^t\one_{\{f\le u\}}=0$, and the supremum of rational
$b>t$ for which $\int_t^b\one_{\{f\le u\}}=0$. These are the two
component endpoints. The tests are measurable by parameterized
integration, so their difference is measurable on the full-measure
set of relevant pairs. Fubini justifies the conventions elsewhere.

The layer-cake identity gives $Df=\int D\one_{E_u}\,du$ as
distributions. For a finite interval union,
$D\one_{E_u}=\sum(\delta_a-\delta_b)$. The localized coarea identity
$|Df|=\int|D\one_{E_u}|\,du$ identifies finite measures, not only
total masses. Adding and subtracting these two identities proves
\eqref{eq:v74-oriented-levels}; see \cite[Chapter 3]{AFP}.
\end{proof}

\begin{theorem}[A level-matched height bound]
\label{thm:v74-level-bound}
For almost every $t$,
\begin{align}
 [f(t)-2A_\delta f(t)]_+
 &\le S_\delta f(t)\notag\\
 &\le\min\{f(t),(Df)^+([t-\delta,t]),
                         (Df)^-([t,t+\delta])\},\label{eq:v74-level-bound}
\end{align}
where $A_\delta f(t)=(2\delta)^{-1}\int_{t-\delta}^{t+\delta}f(s)\,ds$.
Furthermore $S_\delta(cf)=cS_\delta f$ for $c\ge0$, and $S_\delta f$
is nondecreasing as a function of $\delta$. The reserve two is optimal
among bounds of the first form which charge only components shorter
than $\delta$.
\end{theorem}
\begin{proof}
If $t\in I=(a,b)$, then
\[
 |I\cap(t-\delta,t+\delta)|\ge\min\{\delta,b-a\}.
\]
For a level not containing $t$, the contribution to
$\one_{E_u}(t)-\delta^{-1}|E_u\cap(t-\delta,t+\delta)|$ is nonpositive.
For a level containing $t$, the preceding inequality bounds that
contribution by $(1-\ell_u(t)/\delta)_+$. Integration in $u$ and
layer cake give $f-2A_\delta f\le S_\delta f$. The right side is
nonnegative, so this also bounds the positive part.

The bound $S_\delta f\le f$ is immediate. A component contributing
positively has length less than $\delta$; both its entrance and its
exit lie in the corresponding one-sided windows about $t$. Each
contributing weight is at most one. Equation
\eqref{eq:v74-oriented-levels} proves the other two bounds.
The substitution $u=cv$ proves homogeneity for $c>0$, and $c=0$
is immediate. Monotonicity in $\delta$ follows from the integrand.

For optimality of the reserve, take $f=\one_{(0,L)}$ with $L>2\delta$.
Its short-excursion function vanishes. For $0<t<\delta$,
$A_\delta f(t)=(t+\delta)/(2\delta)$. Every reserve $K<2$ leaves
$f-KA_\delta f>0$ on an interval of positive length next to zero.
Thus it cannot replace two in this component-only estimate.
\end{proof}

\begin{proposition}[Sharp integrated identity and length tails]
\label{prop:v74-integrated}
One has
\begin{align}
 \int_\R S_\delta f(t)\,dt
 &=\int_0^\infty\sum_{I\subset E_u}
                         |I|(1-|I|/\delta)_+\,du\notag\\
 &\le\frac{\delta}{8}\|Df\|_{\TV}.\label{eq:v74-integrated}
\end{align}
The constant $1/8$ is sharp. Define
\[
 T_f(t,r)=\int_0^{f(t)}\one_{\{\ell_u(t)<r\}}\,du,
 \qquad
 \Lambda_\beta f(t)=\int_0^{f(t)}\ell_u(t)^{-\beta}\,du.
\]
For $\beta>0$, allowing an infinite moment,
\begin{equation}\label{eq:v74-tail-moment}
 S_\delta f(t)=\frac1\delta\int_0^\delta T_f(t,r)\,dr
 \le\frac{\beta^\beta}{(\beta+1)^{\beta+1}}
                              \delta^\beta\Lambda_\beta f(t).
\end{equation}
\end{proposition}
\begin{proof}
Tonelli gives the identity in \eqref{eq:v74-integrated} by integrating
first on each interval. The bound $x(1-x/\delta)_+\le\delta/4$
and $\int_0^\infty\#\{I\subset E_u\}\,du=\|Df\|_{\TV}/2$
prove its inequality. A pulse of width $\delta/2$ gives equality.
Next $(1-\ell/\delta)_+=\delta^{-1}\int_0^\delta
\one_{\{\ell<r\}}\,dr$, which proves the tail identity.
Finally the maximum of $x^\beta(1-x)$ on $[0,1]$ is
$\beta^\beta/(\beta+1)^{\beta+1}$. Apply this with $x=\ell/\delta$
when $\ell<\delta$, and integrate.
\end{proof}

These bounds are strictly more selective than unpaired Jordan masses.
For $\delta=1$ take
\[
 f=\one_{(-1/2,2)}+\one_{(1/2,3/4)}.
\]
At $t=0$ the incoming positive and outgoing negative Jordan masses
are both one. The levels containing $t$ all have component length
$5/2$, however, so $S_1f(0)=0$. The same strict separation holds
on a neighborhood of zero. The two nearby atoms lie on different
height layers and do not form a short excursion through that roof.
Conversely, for a pulse $f=h\one_{(a,b)}$ with $b-a<\delta$,
\[
 S_\delta f(t)=h(1-(b-a)/\delta)
             =[f(t)-2A_\delta f(t)]_+\quad(a<t<b).
\]
A genuine narrow peak is therefore retained, not canceled. Neither
example is asserted to be a physical Lorentz source. The integrated
bound \eqref{eq:v74-integrated} is an $L^1$ bound only; it is not
converted into a pointwise collision-normalized estimate.
